\section{Kaggle implementation: simulate limited labels}
\label{sec:kaggle-breast-limited-labels}

\begin{kaggleproject}{Breast Cancer Wisconsin --- self-training and uncertainty sampling}
\kagglesource{https://www.kaggle.com/datasets/uciml/breast-cancer-wisconsin-data}{Breast Cancer Wisconsin (Diagnostic) Data Set}
\textbf{Question.} What changes when most training labels are hidden, and how can an active learner prioritize which unknown labels to request next?
\end{kaggleproject}

The outer test split is created first and remains untouched. The remaining development data are divided into a label-acquisition pool and a validation set. The first half hides 70\% of the pool labels and evaluates self-training on validation data. The second half compares uncertainty sampling with a same-budget random-query baseline: both begin from the same small labeled subset, repeatedly fit logistic models, and acquire additional pool labels. Because this is a simulation, the original Kaggle labels act as the oracle that answers each query. Only after the acquisition policy is complete are the final models evaluated on the outer test set once.

\textbf{Before running.} Predict the qualitative pattern you expect from the experiment before you look at the output or plot.

\begin{labmode}
Stage 5: design the label-acquisition evaluation boundary yourself. Development validation may guide the policy; the final test set may not.
\end{labmode}

\textbf{Part A --- create the evaluation boundary and run self-training.}

\lstinputlisting[style=mlpython,firstline=1,lastline=42]{all_experiments/lab18-breast_cancer_limited_labels.py}

\textbf{Part B --- initialize equal label budgets and fit both acquisition policies.}

\lstinputlisting[style=mlpython,firstline=44,lastline=84]{all_experiments/lab18-breast_cancer_limited_labels.py}

\textbf{Part C --- acquire new labels, plot label efficiency, and finish development.}

\lstinputlisting[style=mlpython,firstline=86,lastline=115]{all_experiments/lab18-breast_cancer_limited_labels.py}

\textbf{Part D --- lock the acquisition policy and evaluate the outer test once.}

\lstinputlisting[style=mlpython,firstline=117,lastline=127]{all_experiments/lab18-breast_cancer_limited_labels.py}


\begin{keyidea}
The code is deliberately a simulation of the labeling process. In a real active-learning system, the queried labels would come from a human expert, measurement process, or delayed outcome rather than from a hidden array already available to the program.
\end{keyidea}

\begin{labdeliverable}
Submit the self-training result, an active-learning curve measured on a development validation set, a same-budget random-query baseline, and one final untouched-test evaluation after the labeling policy is locked. Explain why repeatedly checking the test set during label acquisition would turn it into development data.
\end{labdeliverable}
